\lecture{21}
\title{Bayesian Inference I}

\begin{document}

\subsection{Review: Likelihood and Multivariate Distributions}

Suppose we have $n$ independent random variables $X_1, \dots, X_n$,
each with a PDF $f_i(x_i \mid \theta)$ with a dependency
on a parameter $\theta$.
As an example, we might have something like $X_i \sim \mathcal{N}(\theta, i)$.

If we're given the values $X_1, \dots, X_n$,
that lets us compute a likelihood function $L_n(\theta)$
that tells us how likely it would be
to have observed this data given that the parameter had value $\theta$.
(Recall \href{9.html#defining-the-maximum-likelihood-estimator-mle}{Lecture 9}.)
\[ L_n(\theta) := \prod_{i = 1}^n f_i(X_i \mid \theta)
= f(X_1, \dots, X_n \mid \theta). \]
In the above, $f(x_1, \dots, x_n \mid \theta)$ is the joint PDF.
(Recall \href{6.html#multivariate-pdfs}{Lecture 6}.)

\subsection{Priors and Posteriors}

Now suppose we have some \emph{prior} belief
about the distribution of the parameter $\theta$.
As an example, perhaps $\theta$ is a probability
that we believe to be more likely to be high than low.

\textbf{Definition.}
A \textbf{prior} is a PDF $f(\theta)$
that describes the probability distribution
taken on by the parameter $\theta$.

We now read in some data $X_1, \dots, X_n$
that has a dependency on $\theta$.
Bayes' rule tells us the following:
\[ \mathrm{Pr}(\theta \mid X_1, \dots, X_n)
= \frac{
    \mathrm{Pr}(X_1, \dots, X_n \mid \theta) \cdot \mathrm{Pr}(\theta)
}{
    \mathrm{Pr}(X_1, \dots, X_n)
}
= \frac{L_n(\theta) \cdot f(\theta)}{\mathrm{Pr}(X_1, \dots, X_n)}
\propto L_n(\theta) \cdot f(\theta). \]
The LHS is our updated, \emph{posterior} belief
about the distribution of the parameter $\theta$.

\begin{remark}
    Notice how we don't care about
    the exact value of $\mathrm{Pr}(X_1, \dots, X_n)$,
    because it is effectively a normalizing constant
    that works only to make the integral of
    $\mathrm{Pr}(\theta \mid X_1, \dots, X_n)$
    equal to one.
\end{remark}

\textbf{Definition.}
The posterior $f(\theta \mid X_1, \dots, X_n)$
is the probability distribution proportional to
$L_n(\theta) \cdot f(\theta)$.

\subsection{Conjugate Priors}

\textbf{Definition.}
The \textbf{Beta} distribution $\mathrm{Beta}(a, b)$ (where $a, b > 0$)
has PDF $f \colon [0, 1] \to \mathbb{R}^{\geq 0}$
given by $f(p) \propto p^{a - 1} \cdot (1 - p)^{b - 1}$.

\begin{remark}
    For example, the PDF of $\mathrm{Beta}(2, 2)$
    is $f(p) = 6 \cdot p \cdot (1 - p)$.
\end{remark}

\begin{center}
    \frame{\includegraphics[width=0.35\textwidth]{lecture-21/beta-distribution}}
\end{center}

\begin{problem}
    Suppose we sample data i.i.d. from $\mathrm{Ber}(\theta)$,
    and our prior is $\theta \sim \mathrm{Beta}(a, b)$.
    In terms of the data $X_1, \dots, X_n$,
    what is the posterior distribution?
\end{problem}

\begin{solution}
    Suppose $A$ of the samples are one,
    and the remaining $B = n - A$ samples are zero.
    Then we have:
    \begin{align*}
    f(\theta \mid X_1, \dots, X_n) & \propto L_n(\theta) \cdot f(\theta) \\
    & \propto \left[\theta^A \cdot (1 - \theta)^B\right] \cdot
    \left[\theta^{a - 1} \cdot (1 - \theta)^{b - 1}\right] \\
    & = \theta^{(a + A) - 1} \cdot (1 - \theta)^{(b + B) - 1}.
    \end{align*}
    Therefore, the posterior distribution is
    $\mathrm{Beta}(a + A, b + B)$.
\end{solution}

\textbf{Definition.}
If the prior and posterior live in the same family of distributions,
we say that we have a \textbf{conjugate prior}.

\begin{remark}
    In the previous example,
    $\mathrm{Beta}(a, b)$ was a conjugate prior.
    But this fact relies on the distributions of the sample data
    all being $X_i \sim \mathrm{Ber}(\theta)$, rather than, say,
    $X_i \sim \mathcal{N}(\theta, i)$.
\end{remark}

\subsection{Bayes Estimator and Max A Posteriori}

As we read in more data,
our posterior distribution becomes more informative.
In fact, one can show that the standard deviation
of the posterior decreases as $1/\sqrt{n}$
with the number of samples $n$.

\begin{center}
    \frame{\includegraphics[width=0.35\textwidth]{lecture-21/posterior-accuracy}}
\end{center}

Eventually, we might want to summarize the posterior distribution
$f(\theta \mid X_1, \dots, X_n)$ with a single point estimate for $\theta$.

\textbf{Definition.}
The \textbf{Bayes estimator} (or \textbf{mean a posteriori})
is the expected value of $\theta$ under the posterior distribution.

\textbf{Definition.}
The \textbf{max a posteriori (MAP)} is the mode of the posterior;
that is, the value of $\theta$ that maximizes
$f(\theta \mid X_1, \dots, X_n)$.

Generally, the MAP is much easier to compute than the Bayes estimator,
because it doesn't require computing any normalizing factors.

\begin{remark}
    If we have no priors at all---that is,
    our prior $f(\theta)$ is the uniform distribution---then
    the max a posteriori is just the MLE.
    The spirit of Bayesian inference
    is that we must now account for priors.
\end{remark}

\end{document}